\documentclass{article}
\usepackage[utf8]{inputenc}
\usepackage{amsmath}
\usepackage{geometry}
\geometry{margin=2cm}
\begin{document}

\section*{Questão 135 — ENEM 2024 — Ciências da Natureza}

\subsection*{Interpretação da Questão}

A questão apresenta um caso clássico de herança recessiva ligada ao cromossomo X: o daltonismo. Um casal é formado por um homem não daltônico e uma mulher portadora heterozigota do gene para daltonismo. O objetivo é identificar o heredograma que representa corretamente as possibilidades genotípicas e fenotípicas para o bebê desse casal.

\subsection*{Estratégia e Análise}

O primeiro passo é reconhecer que o daltonismo, associado ao X, afeta principalmente homens, pois eles possuem apenas um cromossomo X, e manifestarão a condição se herdarem o gene recessivo. Mulheres precisam de dois alelos recessivos para serem daltônicas e podem ser portadoras se heterozigotas.

Considerando o homem com genótipos XY (X normal) e a mulher heterozigota X\textsuperscript{N} X\textsuperscript{d}, as possibilidades para os filhos são:

\begin{itemize}
  \item Homens: 50\% chances de X\textsuperscript{N}Y (não daltônico), 50\% X\textsuperscript{d}Y (daltônico).
  \item Mulheres: 50\% X\textsuperscript{N}X\textsuperscript{N} (normais), 50\% X\textsuperscript{N}X\textsuperscript{d} (portadoras).
\end{itemize}

Assim, o heredograma deve mostrar essas probabilidades e características.

\subsection*{Conteúdo e Mini Revisão}

O daltonismo é causado por um gene recessivo no cromossomo X. Como os homens só têm um X, qualquer alelo recessivo causa o fenótipo. Em mulheres, apenas homozigotas para o alelo recessivo manifestam a deficiência. Portanto, o padrão da herança gera mais homens daltônicos e mulheres portadoras.

\begin{table}[h!]
\centering
\begin{tabular}{|c|c|c|}
\hline
\textbf{Sexo} & \textbf{Genótipos Possíveis} & \textbf{Fenótipo} \\
\hline
Masculino & X\textsuperscript{N}Y & Normal \\
           & X\textsuperscript{d}Y & Daltônico \\
\hline
Feminino  & X\textsuperscript{N}X\textsuperscript{N} & Normal \\
           & X\textsuperscript{N}X\textsuperscript{d} & Portadora \\
           & X\textsuperscript{d}X\textsuperscript{d} & Daltônica \\
\hline
\end{tabular}
\caption{Genótipos e fenótipos para daltonismo ligado ao cromossomo X}
\end{table}

\subsection*{Resolução e "Pulo do Gato"}

O heredograma da alternativa B representa corretamente a mulher heterozigota portadora do gene recessivo e o homem não afectado. No quadro dos filhos, mostra homens com 50\% chance de ser daltônicos e filhas normais ou portadoras, em concordância com as regras da herança recessiva ligada ao X.

\textbf{Pulo do gato:} lembrar sempre que homens manifestam a condição se carregam o alelo recessivo no X, enquanto mulheres precisam ser homozigotas recessivas para manifestar fenótipo alterado. Essa simplificação facilita a identificação do heredograma correto.

\subsection*{Armadilhas e Alternativas Distratoras}

\begin{itemize}
  \item \textbf{Alternativa A:} Fórmula e gráfico incorretos, desrespeitando a relação genética.
  \item \textbf{Alternativa C:} Confunde herança autossômica com ligada ao X.
  \item \textbf{Alternativa D:} Representação gráfica e fórmula incompatível com o padrão.
  \item \textbf{Alternativa E:} Desconhecimento dos conceitos básicos da genética ligada ao X.
\end{itemize}

\subsection*{Código da Questão}

CN_GENETICA_HEREDOGRAMA_001

\bigskip

\noindent\rule{\textwidth}{0.4pt}

\textbf{Referência bibliográfica:} Griffiths, A. et al. Introdução à genética. Rio de Janeiro: Guanabara Koogan, 2016 (adaptado).

\end{document}